%%=============================================================================
%% Methodologie
%%=============================================================================

\chapter{\IfLanguageName{dutch}{Methodologie}{Methodology}}%
\label{ch:methodologie}

%% TODO: In dit hoofstuk geef je een korte toelichting over hoe je te werk bent
%% gegaan. Verdeel je onderzoek in grote fasen, en licht in elke fase toe wat
%% de doelstelling was, welke deliverables daar uit gekomen zijn, en welke
%% onderzoeksmethoden je daarbij toegepast hebt. Verantwoord waarom je
%% op deze manier te werk gegaan bent.
%% 
%% Voorbeelden van zulke fasen zijn: literatuurstudie, opstellen van een
%% requirements-analyse, opstellen long-list (bij vergelijkende studie),
%% selectie van geschikte tools (bij vergelijkende studie, "short-list"),
%% opzetten testopstelling/PoC, uitvoeren testen en verzamelen
%% van resultaten, analyse van resultaten, ...
%%
%% !!!!! LET OP !!!!!
%%
%% Het is uitdrukkelijk NIET de bedoeling dat je het grootste deel van de corpus
%% van je bachelorproef in dit hoofstuk verwerkt! Dit hoofdstuk is eerder een
%% kort overzicht van je plan van aanpak.
%%
%% Maak voor elke fase (behalve het literatuuronderzoek) een NIEUW HOOFDSTUK aan
%% en geef het een gepaste titel.

Nu we weten dat navigeren op een kerkhof vaak traag of verwarrend is en zelfs met hedendaagse applicaties nog andere zorgen met zich brengt, is het doel concreet om een Proof of Concept op te stellen met als oplossing Augmented Reality(AR) te gebruiken samen met een geoloket als een mobiele applicatie. Het kerkhof van Sint-Gillis-Waas is relatief klein en heeft nog geen publieke applicatie. Hierop kan een geoloket gemaakt worden met een basis AR-overlay die de route aantoont via de mobiele camera. Er kan eventueel echte data gebruikt worden, want de privé applicatie WebBGP gebruikt deze al. Zo niet kan er een kleine dataset aangemaakt worden met mockdata. Deze toepassing wordt enkel gemaakt voor mobiele Ios apparaten. Er wordt geprogrammeerd in Swift. Hierdoor kan er gebruik gemaakt worden van ARKit voor positionele tracking en alle AR zaken. Wanneer de gebruiker zijn gewenste persoon heeft ingevoerd door bepaalde filters (naam, rijksregisternummer, ...) wordt er een bestemming klaargezet. De gebruiker kan snel switchen tussen de camera met de AR functie, of de kaart met een navigatielijn naar de bestemming. We maken dus degelijk gebruik van Augmented Reality door een effectieve pijl en navigatielijn naar de gewenste grafsteen in de omgeving van de gebruiker te brengen die zichtbaar is door de camera van het mobiele apparaat. Als de gebruiker hierdoor rond de omgeving kijkt is het zeer duidelijk waar hij naartoe moet wandelen, ook al is het kerkhof zeer groot of zijn er slechte weersomstandigheden, wat 1 van de factoren was bij een vermoeiend bezoek. \Autocite{Haekkilae2022}. Dit prototype kan getest worden, hiervoor zijn er een aantal succescriteria waarop we kunnen letten: gebruiker maakt X aantal navigatiefouten, ervaart de gebruiker moeite of stress tijdens gebruik, met gebruik van AR is het X aantal \% sneller dan zonder AR, hoe stabiel is de AR, ... Voor dat we het volledige project uitwerken is het maken van een risicoanalyse zeer cruciaal. Belangrijk bij Augmented Reality zijn bepaalde obstakels, hoogteverschillen, het weer of licht die de tracking kan beïnvloeden. Gebruikers die niet vertrouwd zijn met AR moeten nog steeds gebruik kunnen maken van de app zonder enige problemen. Er wordt geen effectieve integratie met de kerkhofadministratie of gemeente verwerkt, enkel een Proof of Concept die eventueel gebruik kan maken van hun data. Als evaluatiemethode vergelijken we verschillende situaties van de bezoeker: geen gebruik van applicatie, gebruik van applicatie zonder AR, gebruik van applicatie met AR. Bij elke test wordt er een technische evaluatie aan gekoppeld om te zien hoe stabiel en hulpvol de Augmented Reality kan zijn. Op Figuur 1 kan de concrete workflow gezien worden. De volgorde is hier cruciaal. De datums zijn een schatting van hoe lang elk deel zou kunnen innemen. Deze zijn dus niet definitief.

%TODO